% ============================================================
%  PREAMBLE — PhD Thesis
%  Topic: Morse-theoretic construction of the Atiyah–Hirzebruch SS
%
%  Robustness: all \newcommand replaced by \providecommand;
%  all \DeclareMathOperator guarded by \let\Foo\relax so that
%  loading this file twice (or alongside an old preamble) is safe.
% ============================================================

% ---- Document class & language -----------------------------------------

\usepackage[english]{babel}

% ---- Page geometry & header/footer ------------------------------------
\usepackage[headheight=18pt]{geometry}
\usepackage{scrlayer-scrpage}
\usepackage{lastpage}

\pagestyle{scrheadings}
\clearpairofpagestyles
\setkomafont{pagehead}{\bfseries}
\ihead{\leftmark}
\ohead{\Large\textbar\enskip\thepage}
\cfoot{}
\automark[section]{section}
\KOMAoptions{headsepline=true}

% ---- Core math packages -----------------------------------------------
\usepackage{amsmath, amsfonts, amsthm}
\usepackage{mathtools}
\usepackage{bbm}
\usepackage{MnSymbol}

% ---- Graphics & figures -----------------------------------------------
\usepackage{graphicx}
\usepackage{float}
\usepackage{wrapfig}
\usepackage{caption}
\usepackage{svg}                % requires Inkscape on PATH

% ---- TikZ & PGF -------------------------------------------------------
\usepackage{tikz}
\usepackage{tikz-cd}
\usepackage{pgfplots}
\pgfplotsset{compat=newest}
\usepgfplotslibrary{colormaps}

\usetikzlibrary{backgrounds, patterns}
\tikzset{>=latex}
\tikzcdset{scale cd/.style={every label/.append style={scale=#1},
		cells={nodes={scale=#1}}}}

% ---- Color & boxes ----------------------------------------------------
\usepackage[dvipsnames]{xcolor}
\usepackage[most]{tcolorbox}
\usepackage{mdframed}
\usepackage{framed}
\usepackage[ntheorem]{empheq}

% ---- Tables -----------------------------------------------------------
\usepackage{array}
\usepackage{tabularx}
\usepackage{longtable}
\usepackage{booktabs}

% ---- Miscellaneous ----------------------------------------------------
\usepackage{pdflscape}
\usepackage[utf8]{inputenc}
\usepackage{fancyvrb}
\usepackage[shortlabels]{enumitem}
\usepackage{todonotes}
\usepackage{adjustbox}
\usepackage{pdfpages}
\usepackage{epstopdf}
\usepackage{etoolbox}
\usepackage{csquotes} 

% ---- Hyperlinks (load late, before cleveref) --------------------------
\usepackage[hidelinks]{hyperref}
\usepackage{cleveref}

% ---- Bibliography -----------------------------------------------------
\usepackage[backend=biber, style=alphabetic]{biblatex}
\addbibresource{sample.bib}

% ============================================================
%  THEOREM ENVIRONMENTS
% ============================================================

\mdfsetup{%
	ntheorem        = false,
	leftmargin      = 10pt,
	everyline       = true,
	splittopskip    = 15pt,
	splitbottomskip = 10pt,
	linewidth       = 1pt
}

\theoremstyle{plain}
\newmdtheoremenv{theorem}{Theorem}[subsection]
\newmdtheoremenv{lemma}[theorem]{Lemma}
\newmdtheoremenv{fact}[theorem]{Fact}

\theoremstyle{definition}
\newtheorem{example}[theorem]{Example}
\newtheorem{examples}[theorem]{Examples}
\newtheorem{definition}[theorem]{Definition}
\newtheorem{defn}[theorem]{Definition}
\newtheorem{corollary}[theorem]{Corollary}
\newtheorem{remark}[theorem]{Remark}
\newtheorem{proposition}[theorem]{Proposition}

% Proof strategy box (left bar style)
\mdfdefinestyle{proofstyle}{%
	leftline         = true,
	rightline        = false,
	topline          = false,
	bottomline       = false,
	linecolor        = gray,
	linewidth        = 3pt,
	innerleftmargin  = 10pt,
	innerrightmargin = 0pt,
	innertopmargin   = 0pt,
	innerbottommargin= 5pt,
	skipabove        = 0pt,
	skipbelow        = 0pt,
	frametitle       = {Proof Strategy:},
	frametitlefont   = \bfseries,
	frametitlerule   = false,
}
\newmdenv[style=proofstyle]{Iprove}

% Numbered diagram box
\newenvironment{fdiagram}{%
	\refstepcounter{equation}%
	\begin{tcolorbox}[
		enhanced, sharp corners, boxrule=1pt,
		colback=white, colframe=black,
		fonttitle=\bfseries, coltitle=black,
		attach boxed title to bottom right={yshift=3mm, xshift=-2mm},
		boxed title style={colback=gray!20, colframe=black,
			sharp corners, boxrule=0.5pt},
		title=Diagram~\theequation
		]%
	}{%
	\end{tcolorbox}%
}

% ============================================================
%  PRETTY SYMBOLS
% ============================================================
\let\uglyepsilon\epsilon  \let\epsilon\varepsilon
\let\uglyphi\phi          \let\phi\varphi

\AtBeginEnvironment{pmatrix}{\renewcommand\arraystretch{1.5}}

% ============================================================
%  MACRO DEFINITIONS
%
%  \providecommand  — silently skips if already defined.
%  \let\Foo\relax before \DeclareMathOperator — resets the name
%    so LaTeX allows redeclaration (safe to do on every run).
% ============================================================

% ---- Number fields ----------------------------------------------------
\providecommand{\R}{\mathbb{R}}
\providecommand{\N}{\mathbb{N}}
\providecommand{\Z}{\mathbb{Z}}
\providecommand{\C}{\mathbb{C}}
\providecommand{\Q}{\mathbb{Q}}
\providecommand{\T}{\mathbb{T}}
\providecommand{\B}{\mathbb{B}}
\providecommand{\PP}{\mathbb{P}}

% ---- Calligraphic letters ---------------------------------------------
\providecommand{\CC}{\mathcal{C}}
\providecommand{\Sp}{\mathcal{S}}
\providecommand{\Chi}{\mathcal{X}}

% ---- Categories -------------------------------------------------------
\let\Set\relax    \DeclareMathOperator{\Set}{\mathbf{Set}}
\let\Ab\relax     \DeclareMathOperator{\Ab}{\mathbf{Ab}}
\let\Ring\relax   \DeclareMathOperator{\Ring}{\mathbf{Ring}}
\let\CDiff\relax  \DeclareMathOperator{\CDiff}{\mathbf{CDiff}}
\DeclareMathOperator{\CGWH}{\mathbf{CGWH}}
\let\Top\relax    \DeclareMathOperator{\Top}{\mathbf{CGWH}}
\let\hTop\relax   \DeclareMathOperator{\hTop}{mathbf{hCGWH}}
\let\TopP\relax   \DeclareMathOperator{\TopP}{\mathbf{TopP}}
\let\TopH\relax   \DeclareMathOperator{\TopH}{\mathbf{TopH}}
\let\hTopH\relax  \DeclareMathOperator{\hTopH}{\mathbf{hTopH}}
\let\TopC\relax   \DeclareMathOperator{\TopC}{\mathbf{CWH}}
\let\hTopC\relax  \DeclareMathOperator{\hTopC}{\mathbf{hCWH}}
\let\TopCPair\relax   \DeclareMathOperator{\TopCPair}{\mathbf{CWH^{2}_{cof}}}
\let\hTopCPair\relax   \DeclareMathOperator{\hTopCPair}{\mathbf{hCWH^{2}_{cof}}}
\let\TopCPoint\relax  \DeclareMathOperator{\TopCPoint}{\mathbf{CWH^*_{cof}}}
\let\hTopCPoint\relax  \DeclareMathOperator{\hTopCPoint}{\mathbf{hCWH^{*}_{cof}}}
\let\hTopP\relax  \DeclareMathOperator{\hTopP}{\mathbf{hCGWH^*}}
\let\absgrp\relax \DeclareMathOperator{\absgrp}{\mathbf{AbSemGrp}}
\let\Cvec\relax   \DeclareMathOperator{\Cvec}{\mathbf{Vect}}
\let\Pmod\relax   \DeclareMathOperator{\Pmod}{\mathbf{Pmod}}
\let\Cmod\relax   \DeclareMathOperator{\Cmod}{\mathbf{Mod}}
\let\ExCoup\relax \DeclareMathOperator{\ExCoup}{\mathbf{ExCoup}}
\let\SpS\relax    \DeclareMathOperator{\SpS}{\mathbf{SpS}}
\let\SCS\relax    \DeclareMathOperator{\SCS}{\mathbf{SCS}}
\let\CICR\relax   \DeclareMathOperator{\CICR}{\mathbf{CICR}}
\renewcommand{\hom}{\mathrm{hom}}
% ---- General math operators -------------------------------------------
\DeclareMathOperator{\Iso}{\mathrm{Iso}}
\DeclareMathOperator{\Hom}{\mathrm{Hom}}
\let\O\relax\DeclareMathOperator{\O}{\mathrm{O}}

\newcommand{\bottg}{\beta} 
\newcommand{\bottmap}[1][-]{\bottg\otimes(#1)}
\let\spec\relax   \DeclareMathOperator{\spec}{\mathrm{spec}}
				  \DeclareMathOperator{\colim}{\mathrm{colim}}
\let\conj\relax   \DeclareMathOperator{\conj}{\mathrm{conj}}
\let\sing\relax   \DeclareMathOperator{\sing}{\mathrm{sing}}
\let\Morse\relax  \DeclareMathOperator{\Morse}{\mathrm{Morse}}
\let\gr\relax     \DeclareMathOperator{\Gr}{\mathrm{Gr}}
\let\supp\relax   \DeclareMathOperator{\supp}{\mathrm{supp}}
\let\cl\relax     \DeclareMathOperator{\cl}{\mathrm{cl}}
\let\inti\relax   \DeclareMathOperator{\inti}{\mathrm{int}}
\let\Tor\relax    \DeclareMathOperator{\Tor}{\mathrm{Tor}}
\let\coker\relax  \DeclareMathOperator{\coker}{\mathrm{coker}}
\let\rank\relax   \DeclareMathOperator{\rank}{\mathrm{rank}}
\let\herm\relax   \DeclareMathOperator{\Herm}{\mathrm{Herm}}
\let\End\relax    \DeclareMathOperator{\End}{\mathrm{End}}
\let\id\relax     \DeclareMathOperator{\id}{\mathrm{id}}
\let\Hess\relax   \DeclareMathOperator{\Hess}{\mathrm{Hess}}
\let\Mat\relax    \DeclareMathOperator{\Mat}{\mathrm{Mat}}
\let\gen\relax    \DeclareMathOperator{\gen}{\mathrm{span}}
\let\Idm\relax    \DeclareMathOperator{\Idm}{\mathbbm{1}}
\let\Fix\relax    \DeclareMathOperator{\Fix}{\mathrm{Fix}}
\let\Aut\relax    \DeclareMathOperator{\Aut}{\mathrm{Aut}}
\let\vect\relax   \DeclareMathOperator{\Vect}{\mathrm{Vect}}
\let\Diff\relax   \DeclareMathOperator{\Diff}{\mathrm{Diff}}

\providecommand{\im}{\mathrm{im}}
\providecommand{\conv}{\mathrm{conv}}
\providecommand{\grad}{\mathrm{grad}}
\providecommand{\Lip}{\mathrm{Lip}}
\providecommand{\diag}{\mathrm{diag}}
\providecommand{\sign}{\mathrm{sign}}
\providecommand{\CW}{\mathrm{\scriptscriptstyle CW}}
\providecommand{\Crit}{\mathrm{Crit}}
\providecommand{\ind}{\mathrm{ind}}
\providecommand{\GL}{\mathrm{GL}}
\providecommand{\G}{\mathrm{G}}
\providecommand{\Ob}{\mathrm{Ob}}
\providecommand{\iso}{\mathrm{iso}}
\providecommand{\BU}{\mathrm{BU}}
\providecommand{\Glue}[1]{\; \bigcup\nolimits_{#1} \;}
\providecommand{\glue}[1]{\;\cup_{#1} \;}

% ---- Differential calculus 	--------------------------------------------
\providecommand*\diff{\mathop{}\!\mathrm{d}}
\providecommand{\dat}[1]{\diff_{#1}}

% ---- Restrictions -----------------------------------------------------
\providecommand\restr[2]{\left.#1\right|_{#2}}

% ---- Delimiters & norms -----------------------------------------------
\let\abs\relax
\DeclarePairedDelimiter\abs{\lvert}{\rvert}
\providecommand{\norm}[1]{\left\lVert#1\right\rVert}

% ---- Algebraic Topology -----------------------------------------------

% ---- Chain complexes ------------
	\providecommand{\ChM}{{}^{\mathcal{M}}C}
	\providecommand{\ChCW}{{}^{\mathcal{CW}}C}
	\providecommand{\ChSing}{{}^{\mathcal{S}}C}
	
% ---- Boundary Maps --------------
	\providecommand{\bM}{{}^{\mathcal{M}}\partial}
	\providecommand{\bCW}{{}^{\mathcal{CW}}\partial}
	\providecommand{\bS}{{}^{\mathcal{S}}\partial}
	
% ---- Coboundary Maps ------------
	\providecommand{\cbM}{{}^{\mathcal{M}}\mathrm{d}}
	\providecommand{\cbCW}{{}^{\mathcal{CW}}\mathrm{d}}
	\providecommand{\cbS}{{}^{\mathcal{S}}\mathrm{d}}
	
% ---- Connecting Homomorphism ----
	\providecommand{\connM}{{}^\mathcal{M}\delta}
	\providecommand{\connCW}{{}^\mathcal{CW}\delta}
	\providecommand{\connS}{{}^\mathcal{S}\delta}

%---- Homology Groups -------------
	\providecommand{\HM}{{}^{\mathcal{M}}H}
	\providecommand{\HCW}{H^{\mathrm{CW}}H}
	\providecommand{\HS}{H^{\mathrm{S}}H}

	\providecommand{\conn}{\delta}





% ---- Morse theory -----------------------------------------------------
\providecommand{\Wu}[1]{W^{\mathrm{u}}(#1)}
\providecommand{\Ws}[1]{W^{\mathrm{s}}(#1)}
\providecommand{\Mflow}[2]{\mathcal{M}(#1,#2)}

% ---- K-theory ---------------------------------------------------------
\let\K\relax  \DeclareMathOperator{\K}{\mathrm{K}}
\providecommand{\KO}{\mathrm{KO}}
\providecommand{\KU}{\mathrm{KU}}
\providecommand{\Kvect}[1]{\mathrm{Vect}(#1)}
\providecommand{\GLR}[1]{\GL(#1,\mathbb{R})}

% ---- Spectral sequences -----------------------------------------------
\providecommand{\Epq}[2]{E^{#1}_{#2}}
\providecommand{\drpq}[1]{d^{#1}}

% ---- Misc helpers -----------------------------------------------------
\providecommand{\euler}{\mathrm{e}}
\providecommand{\ii}{\mathrm{i}}
\providecommand{\transcap}{\text{$\cap\kern-0.4em|\kern0.4em$}}
\providecommand{\Exp}{\mathbb{E}\mathrm{xp}}
\providecommand{\stillOpen}{\color{red}{[\textbf{open}]}}
\providecommand{\Whab}{W_{h^{\beta\alpha}}}
\providecommand{\Lab}{\mathcal{F}(\mathfrak{A}^{\alpha},\mathfrak{A}^{\beta})}

\providecommand*\colvec[3][]{%
	\begin{pmatrix}\ifx\relax#1\relax\else#1\\\fi#2\\#3\end{pmatrix}}

\def\quotient#1#2{%
	\left.\raise1ex\hbox{$#1$}\middle/\lower1ex\hbox{$#2$}\right.}